\documentclass[../../main.tex]{subfiles}
\begin{document}

{\bf [解答]}

% TODO: 3d-tikzでのD_1,D_2の様子 (スケッチ上図の完全再現)
% TODO: 3d-tikzでのD_1,D_2の様子 (回転軸補正 & y=s 切断平面追加版)
\begin{figure}[htb]
  \centering
  \begin{tikzpicture}[
      x={(-0.85cm, -0.45cm)},
      y={(1.35cm, -0.15cm)},
      z={(0cm, 1.65cm)},
      scale=1.45, >=stealth, every node/.style={font=\small}
    ]

    % パラメータ設定
    \def\a{1.3}       % 円盤の高さ a
    \def\sVal{0.45}   % 切断面 y = s (0 < s < 1)
    \pgfmathsetmacro{\cVal}{sqrt(1.0 - \sVal*\sVal)} % 弦の半長 c = sqrt(1 - s^2) ≈ 0.893

    % ============================================================
    % 1. 座標軸 (奥側)
    % ============================================================
    \draw[thick, gray!70!black] (0, 0, -2.1) -- (0, 0, 0);
    \draw[dashed, gray!50] (0, -1.8, 0) -- (0, 0, 0);
    \draw[->, thick, gray!70!black] (0, 0, 0) -- (0, 2.3, 0) node[below right] {$y$};
    \draw[->, thick, gray!70!black] (0, 0, 0) -- (0, 0, 2.3) node[above] {$z$};
    \node[below left=2pt] at (0, 0, 0) {$O$};

    % ============================================================
    % 2. 円柱の骨格母線 (破線)
    % ============================================================
    \draw[dashed, gray!45, line width=0.8pt] (1, 0, -\a) -- (1, 0, \a);
    \draw[dashed, gray!45, line width=0.8pt] (-1, 0, -\a) -- (-1, 0, \a);
    \draw[dashed, gray!45, line width=0.8pt] (0, 1, -\a) -- (0, 1, \a);
    \draw[dashed, gray!45, line width=0.8pt] (0, -1, -\a) -- (0, -1, \a);

    % ============================================================
    % 3. 平面 y = s (半透明の薄いオレンジ塗り)
    % ============================================================
    \fill[orange!20, opacity=0.55]
    (-1.25, \sVal, -\a - 0.25) -- (1.25, \sVal, -\a - 0.25)
    -- (1.25, \sVal, \a + 0.35) -- (-1.25, \sVal, \a + 0.35) -- cycle;
    \draw[dashed, orange!80!black, line width=0.8pt]
    (-1.25, \sVal, -\a - 0.25) -- (1.25, \sVal, -\a - 0.25)
    -- (1.25, \sVal, \a + 0.35) -- (-1.25, \sVal, \a + 0.35) -- cycle;
    \node[orange!90!black, font=\footnotesize\bfseries, above]
    at (1.1, \sVal, \a + 0.35) {平面 $y = s$};

    % y 軸上の目盛り s
    \node[circle, fill, inner sep=1.1pt] at (0, \sVal, 0) {}; \node[below=2pt, font=\scriptsize] at (0, \sVal, 0) {$s$};

    % ============================================================
    % 4. 下の円盤 D_2: z = -a
    % ============================================================
    \fill[blue!10, opacity=0.6, domain=0:360, samples=60, variable=\t]
    plot ({cos(\t)}, {sin(\t)}, {-\a});
    \draw[line width=1.3pt, blue!85!black, domain=0:360, samples=60, variable=\t]
    plot ({cos(\t)}, {sin(\t)}, {-\a});
    \node[circle, fill, inner sep=1.2pt] at (0, 0, -\a) {}; \node[below right=1pt, font=\scriptsize] at (0, 0, -\a) {$-a$};
    \node[right=2pt, font=\footnotesize\bfseries, blue!90!black] at (0.8, 0.7, -\a) {$D_2$};

    % D_2 と平面 y = s の交線（太い赤線）
    \draw[line width=2.0pt, red!85!black]
    (-\cVal, \sVal, -\a) -- (\cVal, \sVal, -\a);
    \node[circle, fill=red!85!black, inner sep=1.2pt] at (-\cVal, \sVal, -\a) {};
    \node[circle, fill=red!85!black, inner sep=1.2pt] at (\cVal, \sVal, -\a) {};

    % ============================================================
    % 5. 上の円盤 D_1: z = a
    % ============================================================
    \fill[blue!10, opacity=0.6, domain=0:360, samples=60, variable=\t]
    plot ({cos(\t)}, {sin(\t)}, {\a});
    \draw[line width=1.3pt, blue!85!black, domain=0:360, samples=60, variable=\t]
    plot ({cos(\t)}, {sin(\t)}, {\a});
    \node[circle, fill, inner sep=1.2pt] at (0, 0, \a) {}; \node[above left=1pt, font=\scriptsize] at (0, 0, \a) {$a$};
    \node[right=2pt, font=\footnotesize\bfseries, blue!90!black] at (0.8, 0.7, \a) {$D_1$};

    % D_1 と平面 y = s の交線（太い赤線）
    \draw[line width=2.0pt, red!85!black]
    (-\cVal, \sVal, \a) -- (\cVal, \sVal, \a);
    \node[circle, fill=red!85!black, inner sep=1.2pt] at (-\cVal, \sVal, \a) {};
    \node[circle, fill=red!85!black, inner sep=1.2pt] at (\cVal, \sVal, \a) {};
    \node[red!90!black, font=\scriptsize\bfseries, right=2pt]
    at (\cVal, \sVal, \a) {交線 $[-c, c]$};

    % ============================================================
    % 6. 座標軸 (手前側 x 軸)
    % ============================================================
    \draw[->, thick, gray!70!black] (0, 0, 0) -- (2.0, 0, 0) node[below left] {$x$};

    % ============================================================
    % 7. y 軸まわりの回転矢印 (y 軸の先端付近を z軸から x軸方向へ回す)
    % ============================================================
    % y 軸上の位置 y = 2.0 で、z 軸側 (t = 80°) から x 軸側 (t = 0°) へ回す
    \def\yArrowPos{2.05}
    \def\rRot{0.45}

    \draw[->, line width=1.4pt, red!85!black, domain=85:5, samples=30, smooth, variable=\t]
    plot ({\rRot*sin(\t)}, \yArrowPos, {\rRot*cos(\t)});

    % 回転軸のラベル (y 軸の矢印の先に配置)
    \node[red!90!black, font=\footnotesize\bfseries, above right=2pt]
    at ({0.4*sin(50)}, \yArrowPos, {0.4*cos(50)}) {回転軸 $y$};

  \end{tikzpicture}
  \caption{空間内の2つの円盤 $D_1, D_2$，平面 $y=s$ による切断線分（赤線），および $y$ 軸まわりの回転}
  \label{fig:two_disks_3d_overview_fixed}
\end{figure}

与えられた円盤
\begin{align}
  \begin{cases}
    D_1: x^2+y^2 \le 1, z = a \\
    D_2: x^2+y^2 \le 1, z = -a
  \end{cases} \label{eq:1}
\end{align}
について考える．以下簡単のため$c=\cos\theta,s=\sin\theta$と書く．

\paragraph{(1)}

% TODO: y = sでの2円盤の切断面（それぞれ線分）
\begin{figure}[htb]
  \centering
  \begin{tikzpicture}[x=2.2cm, y=2.2cm, >=stealth, every node/.style={font=\small}]

    % ============================================================
    % 1. 定数設定 (スケッチの比率に合わせた幾何パラメータ)
    % a = 1.0 (円盤の高さ)
    % c = 0.9 (y = s での弦の半長 c = cos theta)
    % 外半径 R = sqrt(a^2 + c^2) = sqrt(1.0 + 0.81) = 1.345
    % ============================================================
    \def\a{1.0}
    \def\cVal{0.9}
    \pgfmathsetmacro{\RVal}{sqrt(\a*\a + \cVal*\cVal)} % ≈ 1.345

    % ============================================================
    % 3. 座標軸 (切断面 y = s 内の xz 座標系)
    % ============================================================
    \draw[->, thick, gray!80!black] (-1.6, 0) -- (1.75, 0) node[right] {$x$};
    \draw[->, thick, gray!80!black] (0, -1.6) -- (0, 1.7) node[above] {$z$};
    \node[below left=2pt] at (0, 0) {$O$};
    \node[above right=1pt, font=\footnotesize\bfseries, gray!70!black] at (0.9, 1.45) {$y = s$};

    % ============================================================
    % 4. 境界線の描画
    % ============================================================

    % 円盤の切断線分 [-c, c] (z = a および z = -a)
    \draw[line width=1.3pt, gray!30!black] (-\cVal, \a) -- (\cVal, \a);
    \draw[line width=1.3pt, gray!30!black] (-\cVal, -\a) -- (\cVal, -\a);

    % 左側の境界垂直線 x = -c (破線)
    \draw[dashed, gray!60, line width=0.8pt] (-\cVal, -\a) -- (-\cVal, \a);
    \draw[dashed, gray!60, line width=0.8pt] (\cVal, -\a) -- (\cVal, \a);

    % ============================================================
    % 5. 動径線・補助破線 (スケッチ通り)
    % ============================================================
    \draw[dashed, gray!45] (0, 0) -- (\cVal, \a);
    \draw[dashed, gray!45] (0, 0) -- (\cVal, -\a);
    \draw[dashed, gray!45] (0, 0) -- (-\cVal, \a);
    \draw[dashed, gray!45] (0, 0) -- (-\cVal, -\a);

    % ============================================================
    % 6. 目盛り・寸法ラベル (スケッチ完全準拠)
    % ============================================================
    % 特徴点のプロット
    \node[circle, fill, inner sep=1.2pt] at (\cVal, \a) {};
    \node[circle, fill, inner sep=1.2pt] at (\cVal, -\a) {};
    \node[circle, fill, inner sep=1.2pt] at (-\cVal, \a) {};
    \node[circle, fill, inner sep=1.2pt] at (-\cVal, -\a) {};

    \node[above=2pt, font=\scriptsize] at (0, \a) {$a$};
    \node[below=2pt, font=\scriptsize] at (0, -\a) {$-a$};
    \node[below=2pt, font=\scriptsize] at (\cVal, 0) {$c$};
    \node[below=2pt, font=\scriptsize] at (-\cVal, 0) {$-c$};

  \end{tikzpicture}
  \caption{平面 $y=s$ における円盤の切断線分}
\end{figure}

% TODO: y=sでの切断面 (スケッチ下図の完全再現)
\begin{figure}[htb]
  \centering
  \begin{tikzpicture}[x=2.2cm, y=2.2cm, >=stealth, every node/.style={font=\small}]

    % ============================================================
    % 1. 定数設定 (スケッチの比率に合わせた幾何パラメータ)
    % a = 1.0 (円盤の高さ)
    % c = 0.9 (y = s での弦の半長 c = cos theta)
    % 外半径 R = sqrt(a^2 + c^2) = sqrt(1.0 + 0.81) = 1.345
    % ============================================================
    \def\a{1.0}
    \def\cVal{0.9}
    \pgfmathsetmacro{\RVal}{sqrt(\a*\a + \cVal*\cVal)} % ≈ 1.345

    % ============================================================
    % 2. 領域の塗りつぶし (ハッチング)
    % ============================================================
    % 右側 (x >= 0): 半円環 S(theta) (斜線ハッチング)
    \fill[pattern=north east lines, pattern color=blue!75, opacity=0.7]
    (0, \a) arc (90:-90:\a) -- (0, -\RVal) arc (-90:90:\RVal) -- cycle;

    % 左側 (x <= 0): T(theta) (手書きスケッチ通りの縦線ハッチング)
    % 上部: y = s における D1 の線分 [-c, 0] と外側円弧の間
    \fill[pattern=vertical lines, pattern color=teal!80, opacity=0.7]
    (-\cVal, \a) -- (0, \a) arc (90:{180 - asin(\a/\RVal)}:\RVal) -- cycle;
    % 下部: D2 の線分 [-c, 0] と外側円弧の間
    \fill[pattern=vertical lines, pattern color=teal!80, opacity=0.7]
    (-\cVal, -\a) -- (0, -\a) arc (-90:{-180 + asin(\a/\RVal)}:\RVal) -- cycle;

    % ============================================================
    % 3. 座標軸 (切断面 y = s 内の xz 座標系)
    % ============================================================
    \draw[->, thick, gray!80!black] (-1.6, 0) -- (1.75, 0) node[right] {$x$};
    \draw[->, thick, gray!80!black] (0, -1.6) -- (0, 1.7) node[above] {$z$};
    \node[below left=2pt] at (0, 0) {$O$};
    \node[above right=1pt, font=\footnotesize\bfseries, gray!70!black] at (0.9, 1.45) {$y = s$};

    % ============================================================
    % 4. 境界線の描画
    % ============================================================
    % 右半面の外円弧 (半径 R = sqrt(a^2+c^2))
    \draw[line width=1.4pt, blue!85!black] (0, \RVal) arc (90:-90:\RVal);
    % 右半面の内円弧 (半径 a)
    \draw[line width=1.3pt, blue!85!black] (0, \a) arc (90:-90:\a);

    % 左半面の外円弧
    \draw[line width=1.4pt, teal!85!black]
    (0, \RVal) arc (90:{180 - asin(\a/\RVal)}:\RVal);
    \draw[line width=1.4pt, teal!85!black]
    (0, -\RVal) arc (-90:{-180 + asin(\a/\RVal)}:\RVal);

    % 円盤の切断線分 [-c, c] (z = a および z = -a)
    \draw[line width=1.3pt, gray!30!black] (-\cVal, \a) -- (\cVal, \a);
    \draw[line width=1.3pt, gray!30!black] (-\cVal, -\a) -- (\cVal, -\a);

    % 左側の境界垂直線 x = -c (破線)
    \draw[dashed, gray!60, line width=0.8pt] (-\cVal, -\a) -- (-\cVal, \a);

    % ============================================================
    % 5. 動径線・補助破線 (スケッチ通り)
    % ============================================================
    \draw[dashed, gray!45] (0, 0) -- (\cVal, \a);
    \draw[dashed, gray!45] (0, 0) -- (\cVal, -\a);
    \draw[dashed, gray!45] (0, 0) -- (-\cVal, \a);
    \draw[dashed, gray!45] (0, 0) -- (-\cVal, -\a);

    % ============================================================
    % 6. 目盛り・寸法ラベル (スケッチ完全準拠)
    % ============================================================
    % 特徴点のプロット
    \node[circle, fill, inner sep=1.2pt] at (\cVal, \a) {};
    \node[circle, fill, inner sep=1.2pt] at (\cVal, -\a) {};
    \node[circle, fill, inner sep=1.2pt] at (-\cVal, \a) {};
    \node[circle, fill, inner sep=1.2pt] at (-\cVal, -\a) {};

    \node[above=2pt, font=\scriptsize] at (0, \a) {$a$};
    \node[below=2pt, font=\scriptsize] at (0, -\a) {$-a$};
    \node[below=2pt, font=\scriptsize] at (\cVal, 0) {$c$};
    \node[below=2pt, font=\scriptsize] at (-\cVal, 0) {$-c$};

    % 外半径・内半径の引き出し線 (スケッチの指示矢印)
    \draw[<-, gray!30!black, line width=0.7pt]
    ({0.5*\RVal*cos(20)}, {0.5*\RVal*sin(20)}) ++({0.5*\RVal*cos(20)}, {0.5*\RVal*sin(20)})
    to[out=30, in=180] (1.35, 0.45) node[right, font=\footnotesize\bfseries] {$\sqrt{a^2+c^2}$};

    \draw[<-, gray!30!black, line width=0.7pt]
    (\a, 0) to[out=-30, in=180] (1.2, -0.25) node[right, font=\footnotesize\bfseries] {$a$};

    % 領域ラベル
    \node[blue!90!black, font=\normalsize\bfseries] at (0.9, -0.7) {$S(\theta)$};

  \end{tikzpicture}
  \caption{平面 $y=s$ における切断面（$x \ge 0$ の半円環 $S(\theta)$）}
  \label{fig:cross_section_ys_Stheta}
\end{figure}

対称性から$y \ge 0$で考えよう．
$y = s \ (0 \le \theta \le \pi/2)$で切断すると\cref{eq:1}は
\begin{align}
  \begin{cases}
    D_1: -c \le x \le c, z = a \\
    D_2: -c \le x \le c, z = -a
  \end{cases} \label{eq:2}
\end{align}
という二本の線分であり，
この部分を $y$ 軸まわりに回した時の概形は上図のようになる．

特に $x \ge 0$ をみたす部分の面積 $S(\theta)$ は
半径$\sqrt{a^2+c^2}$と$a$の二つの円の間のドーナツ型の部分の面積だから
\begin{align}
  S(\theta) = \dfrac{1}{2}\pi \left[ (a^2+c^2) - a^2 \right] = \dfrac{1}{2}\pi c^2
\end{align}
だから，対称性から求める体積$W(a)$は
\begin{align}
  \dfrac{1}{2} W(a)
  &= \int_{0}^{s} S(\theta) ds \\
  &= \int_{0}^{\pi/2} \dfrac{1}{2}\pi c^2 \dfrac{ds}{d\theta} d\theta \\
  &= \dfrac{1}{2}\pi \int_{0}^{\pi/2} c^3 d\theta \\
  &= \dfrac{1}{2}\pi \int_{0}^{\pi/2} \left( \dfrac{3}{4}c + \dfrac{1}{4}c \cos 3\theta \right) d\theta \\
  &= \dfrac{1}{2}\pi \left[ \dfrac{3}{4}s + \dfrac{1}{12}\sin 3\theta \right]_{0}^{\pi/2} \\
  &= \dfrac{1}{2}\left(\dfrac{3}{4} - \dfrac{1}{12}\right)\pi = \dfrac{1}{3}\pi
\end{align}
となって，
\begin{align}
  W(a) = \dfrac{2}{3}\pi \label{eq:3}
\end{align}
である．

\paragraph{(2)}

% TODO: 再びy=sでの切断面，x<0にフォーカス
% TODO: 再びy=sでの切断面，x<0にフォーカス
\begin{figure}[htb]
  \centering
  \begin{tikzpicture}[x=2.8cm, y=2.4cm, >=stealth, every node/.style={font=\small}]

    % パラメータ
    \def\a{1.0}
    \def\cVal{0.9}
    \pgfmathsetmacro{\RVal}{sqrt(\a*\a + \cVal*\cVal)} % ≈ 1.345

    % ============================================================
    % 1. 領域 T(theta) の縦線ハッチング (手書きスケッチ通り)
    % 上部: z = a から z = sqrt(R^2 - x^2)
    % 下部: z = -sqrt(R^2 - x^2) から z = -a
    % ============================================================
    \fill[pattern=vertical lines, pattern color=teal!85!black, opacity=0.75]
    (-\cVal, \a) -- (0, \a) -- (0, \RVal)
    arc (90:{180 - asin(\a/\RVal)}:\RVal) -- cycle;

    \fill[pattern=vertical lines, pattern color=teal!85!black, opacity=0.75]
    (-\cVal, -\a) -- (0, -\a) -- (0, -\RVal)
    arc (-90:{-180 + asin(\a/\RVal)}:\RVal) -- cycle;

    % ============================================================
    % 2. 座標軸
    % ============================================================
    \draw[->, thick, gray!80!black] (-1.55, 0) -- (0.4, 0) node[right] {$x$};
    \draw[->, thick, gray!80!black] (0, -1.6) -- (0, 1.6) node[above] {$z$};
    \node[below right=2pt] at (0, 0) {$O$};

    % ============================================================
    % 3. 境界線
    % ============================================================
    % 外円弧 z = ±sqrt(a^2 + c^2 - x^2)
    \draw[line width=1.5pt, teal!85!black, domain={-\cVal}:0, samples=50, smooth, variable=\x]
    plot ({\x}, {sqrt(\RVal*\RVal - \x*\x)});
    \draw[line width=1.5pt, teal!85!black, domain={-\cVal}:0, samples=50, smooth, variable=\x]
    plot ({\x}, {-sqrt(\RVal*\RVal - \x*\x)});

    \node[teal!90!black, font=\footnotesize\bfseries, above left] at (-0.45, 1.25)
    {$z = \sqrt{a^2+c^2-x^2}$};

    % 円盤の線分 z = ±a
    \draw[line width=1.4pt, gray!30!black] (-\cVal, \a) -- (0, \a);
    \draw[line width=1.4pt, gray!30!black] (-\cVal, -\a) -- (0, -\a);
    \node[above=1pt, font=\scriptsize] at (-0.45, \a) {$z = a$};
    \node[below=1pt, font=\scriptsize] at (-0.45, -\a) {$z = -a$};

    % 端点 x = -c の垂直破線
    \draw[dashed, gray!60, line width=0.9pt] (-\cVal, -\a) -- (-\cVal, \a);

    % ============================================================
    % 4. 目盛りと点プロット
    % ============================================================
    \node[circle, fill, inner sep=1.3pt] at (-\cVal, \a) {}; \node[left=2pt, font=\scriptsize] at (-\cVal, \a) {$( -c, a )$};
    \node[circle, fill, inner sep=1.3pt] at (-\cVal, -\a) {}; \node[left=2pt, font=\scriptsize] at (-\cVal, -\a) {$( -c, -a )$};
    \node[circle, fill, inner sep=1.2pt] at (0, \a) {}; \node[right=2pt, font=\scriptsize] at (0, \a) {$a$};
    \node[circle, fill, inner sep=1.2pt] at (0, -\a) {}; \node[right=2pt, font=\scriptsize] at (0, -\a) {$-a$};
    \node[circle, fill, inner sep=1.2pt] at (0, \RVal) {}; \node[right=2pt, font=\scriptsize] at (0, \RVal) {$\sqrt{a^2+c^2}$};

    \node[below=2pt, font=\scriptsize] at (-\cVal, 0) {$-c$};

    % 積分微小幅 dx のインジケータ
    \draw[<->, red!85!black, line width=1.0pt] (-0.5, \a) -- (-0.5, {sqrt(\RVal*\RVal - 0.25)})
    node[midway, left=1pt, font=\scriptsize\bfseries] {高さ};

    % 領域ラベル T(theta)
    \node[teal!95!black, font=\large\bfseries] at (-0.5, 0.55) {$\frac{1}{2}T(\theta)$};
    \node[teal!95!black, font=\large\bfseries] at (-0.5, -0.55) {$\frac{1}{2}T(\theta)$};

  \end{tikzpicture}
  \caption{$x \le 0$ における切断面の領域 $T(\theta)$（上下の縦線部）}
  \label{fig:cross_section_Ttheta_x_neg}
\end{figure}

$E$ のうち $x \le 0$ のものの $y=s$ における切断面の面積 $T(\theta)$ は，
\begin{align}
  T(\theta)
  &= 2 \int_{-c}^{0} \left[ \sqrt{a^2+c^2-x^2} - a \right] dx \\
  &= 2 \int_{0}^{c} \sqrt{a^2+c^2-x^2} dx - 2ac \label{eq:4}
\end{align}
であり，これを$y$について積分すれば$V(a)$のうち$x \le 0$の部分になるから
\begin{align}
  V(a) = W(a) + 2 \int_{0}^{1} T(\theta) ds \label{eq:5}
\end{align}
となる．以下$a\to\infty$の極限を考えるので，\cref{eq:5}の第二項の
$a\to\infty$での極限値を考えれば良い．

ここで，$[0,c]$ において，$0 \le x^2 \le c^2$ だから，
\begin{align}
  \sqrt{a^2} \le \sqrt{a^2+c^2-x^2} \le \sqrt{a^2+c^2}
\end{align}
同区間$[0,c]$で積分して，$a>0$ から$\sqrt{a^2}=a$なので
\begin{align}
  & \int_{0}^{c} a dx \le \int_{0}^{c} \sqrt{a^2+c^2-x^2} dx \le \int_{0}^{c} \sqrt{a^2+c^2} dx \\
  \therefore
  & ac \le \int_{0}^{c} \sqrt{a^2+c^2-x^2} dx \le c\sqrt{a^2+c^2}
\end{align}
だから，\cref{eq:4}より
\begin{align}
  0 \le T(\theta) \le 2c \left[ \sqrt{a^2+c^2} - a \right] \label{eq:6}
\end{align}
である．これを$[0,1]$ で積分して，
\begin{align}
  0 \le \int_{0}^{1} T(\theta) dy \le 2 \int_{0}^{1} c \left(\sqrt{a^2+c^2} - a \right) c dy \label{eq:7}
\end{align}
である．$y=s$より$\dfrac{dy}{d\theta}=c$だから置換積分によって
\begin{align}
  0 \le \int_{0}^{1} T(\theta) dy \le 2 \int_{0}^{\pi/2} c^2\left( \sqrt{a^2+c^2} - a \right) d\theta \le 2(\sqrt{a^2+1} - a) \int_{0}^{pi/2} c^2 d\theta
\end{align}
である．ここで最後の不等号は$\sqrt{a^2+c^2} - a \le \sqrt{a^2+1} -a$より従う．
再右辺に$\dfrac{\sqrt{a^2+1} + a}{\sqrt{a^2+1} + a}$をかけて
\begin{align}
  0 \le \int_{0}^{1} T(\theta) dy \le 2\dfrac{1}{\sqrt{a^2+1} + a} \int_{0}^{pi/2} c^2 d\theta
\end{align}
である．$a\to\infty$で右辺は$0$に収束するから挟み撃ちの定理より
\begin{align}
  \lim_{a\to\infty} \int_{0}^{1} T(\theta) dy = 0
\end{align}
であり，\cref{eq:5}に代入して(1)の結果と合わせて，求める極限値は
\begin{align}
  \lim_{a\to\infty} V(a) = \lim_{a\to\infty} W(a) = \dfrac{2}{3}\pi
\end{align}
である．

\end{document}
